\section{System Overview}

Our final system is a self-contained four-stage pipeline.  First, a sparse
candidate generator retrieves top-500 evidence with BM25-style word, character
n-gram, structured cue, and PRF/query-expansion signals, fused by reciprocal
rank fusion.  Second, a MiniLM embedding scorer and shallow factual features
rerank the candidate pool to top-64.  Third, a bounded cross-encoder pass is
combined with embedding and source-rank signals to select the final top-3
evidence.  Finally, claim labels are predicted with an enhanced Context20
TF-IDF one-vs-rest logistic regression classifier.

\section{Enhanced Context20 Classifier}

The original Context20 classifier used a plain concatenation format:
\texttt{CLAIM} followed by twenty \texttt{EVIDENCE} lines.  This loses two
important signals.  First, rank 1 and rank 20 evidence become almost equivalent
to TF-IDF.  Second, factual matching cues used by the top-64 reranker, such as
word overlap, number overlap, and negation compatibility, are not explicitly
available to the classifier.

We therefore use an enhanced Context20 representation.  Evidence ranked 1--5 is
prefixed as \texttt{TOP\_EVIDENCE}, evidence ranked 6--20 as
\texttt{MID\_EVIDENCE}, and the top five evidence passages are repeated once to
increase their TF-IDF weight.  Each evidence line also receives compact cue
tokens for word-overlap bucket, number match/mismatch, and negation
match/mismatch.  These tokens are treated only as input features; the classifier
still learns all label decisions through logistic regression.  The final
classifier uses word 1--2 gram TF-IDF with 60,000 features, balanced
one-vs-rest logistic regression, and $C=0.125$.

\section{Development Results}

\begin{table}[h]
\centering
\begin{tabular}{lrr}
\hline
Stage & Metric & Value \\
\hline
Sparse candidate top-500 & Macro recall & 0.6754 \\
Embedding + hand top-64 & Macro recall & 0.5784 \\
CE + embedding/source top-3 & Macro recall & 0.2355 \\
Final evidence top-3 & Evidence F & 0.2021 \\
Enhanced Context20 classifier & Accuracy & 0.5195 \\
Enhanced Context20 classifier & Macro-F1 & 0.4841 \\
\hline
\end{tabular}
\caption{Development-set results for the final self-contained pipeline.}
\end{table}

\section{Hyperparameter Sensitivity}

The classifier was most sensitive to the logistic regression regularization
strength.  Stronger regularization improved both accuracy and macro-F1.

\begin{table}[h]
\centering
\begin{tabular}{rrr}
\hline
$C$ & Accuracy & Macro-F1 \\
\hline
0.125 & 0.5195 & 0.4841 \\
0.25  & 0.4935 & 0.4544 \\
0.5   & 0.4740 & 0.4263 \\
1.0   & 0.4675 & 0.4184 \\
2.0   & 0.4416 & 0.3825 \\
4.0   & 0.4545 & 0.3913 \\
8.0   & 0.4481 & 0.3925 \\
\hline
\end{tabular}
\caption{Development sensitivity to logistic regression regularization.}
\end{table}

\begin{table}[h]
\centering
\begin{tabular}{lrr}
\hline
Context format & Accuracy & Macro-F1 \\
\hline
Plain Context20 & 0.4351 & 0.3762 \\
Rank-aware prefix & 0.4351 & 0.3749 \\
Top5 weighted & 0.4416 & 0.3938 \\
Rank weighted & 0.4351 & 0.3882 \\
Enhanced Context20, $C=0.25$ & 0.4935 & 0.4544 \\
Enhanced Context20, $C=0.125$ & 0.5195 & 0.4841 \\
\hline
\end{tabular}
\caption{Development ablation for context formatting.}
\end{table}

\section{Model Selection Caveat}

The final $C=0.125$ setting was selected from a development diagnostic
sensitivity analysis, not from a purely train-only fold selection.  However,
train-fold diagnostics supported the main design choice: the enhanced context
representation improved macro-F1 compared with the plain Context20 input,
indicating that rank and factual cue tokens provide useful signal rather than
only fitting one development split.
